\documentclass[10pt,a4paper]{amsart}
\usepackage[margin=19mm]{geometry}
\usepackage{amsmath,amssymb,mathtools}
\usepackage[scr=rsfs,cal=euler]{mathalfa}
\usepackage{xfrac}
\usepackage[unicode,hidelinks,bookmarks=false]{hyperref}
\newcommand{\ie}{i.e.\ }
\newcommand{\cf}{cf.\ }
\newcommand{\resp}{respectively\ }
\newcommand{\Subsection}{\subsection}
\providecommand{\nobreakdash}{\nobreak\hbox{-}}
\setlength{\emergencystretch}{2em}
\multlinegap0pt
\usepackage{mathtools}
\usepackage{amssymb}
\usepackage{enumerate}
\usepackage{float}
\usepackage{tikz}
\usepackage{dsfont}
\newtheorem{theorem}{Theorem}
\title{Strong majority colorings of graphs}
\author{Rafa{\l} Kalinowski, Mateusz Kamyczura, Monika Pil\'sniak, Mariusz Wo\'zniak}
\date{Source: arXiv:2605.23828v1 (CC BY 4.0)}
\begin{document}
\maketitle

\noindent\emph{Source and licence.} Strong majority colorings of graphs. Version arXiv:2605.23828v1 (CC BY 4.0), distributed by arXiv under the Creative Commons Attribution 4.0 International licence, http://creativecommons.org/licenses/by/4.0/, as recorded in the arXiv licence metadata for this exact version and shown on the abstract page https://arxiv.org/abs/2605.23828v1. Full licence notice: \texttt{licenses/arxiv-2605.23828-CC-BY-4.0.txt}.\par\smallskip

\noindent\textit{Editorial scope.} Counted unit: the article's theorem col (source0.tex lines 334--336). The article's complete original proof of it is delivered with it (source0.tex lines 338--365, one \texttt{proof} environment of 28 lines). No mathematics is written, repaired or re-derived: the statement and the proof are the article's own lines, in the article's own notation, and no retained line is substituted or re-flowed.  No further source-local context is needed: every cross-reference of every retained line resolves inside the two delivered spans.  Nothing was omitted from any retained span, so the package carries no omission record.  No retained line contains a citation, so the wrapper carries no bibliography and no bibliography entry.  No retained line contains a figure environment, an image-inclusion command or drawing code, so no image is delivered and the package declares no asset dependency.  Editorial formatting only, in addition: an AMS \texttt{amsart} preamble that restates the article's own declarations and macros used by the retained lines, namely the environments \texttt{theorem} on separate counters, together with the standard packages those lines need.  The retained environments print in this document as Theorem 1 in order of appearance, whereas the article prints the same items with the numbering of its own full document.  The complete native source of the article as distributed in the arXiv e-print for this exact version is bundled unchanged as \texttt{sources/201238/source0.tex}.
\begin{theorem}\label{edge col}
For every admissible graph $G$, $${\mathrm Maj'}(G)\le 8.$$    
\end{theorem}

\begin{proof} Let $G=(V,E)$ be an admissible graph, and let $C$ be a set of eight colors. We want to show that there is a strong majority edge-coloring of $G$ with  colors of $C$. We may assume that $G$ is connected; otherwise we can argue component-wise.

We construct a graph $G^*$ using the following operation of splitting vertices. If $\Delta(G)\le 3$, then we put  $G^*=G$.  Otherwise, for every vertex $v\in V$  of degree greater than 3, we partition its neighborhood $N(v)$ into subsets $N_1(v),\ldots,N_s(v)$ such that $|N_i(v)|\in \{2,3\}, i=1,\ldots,s$, and the number of $i$'s with $|N_i(v)|=2$ is at most two. Next, we substitute the vertex $v$ by $s$ new vertices $v_1,\ldots,v_s$, and insert an edge between $v_i$ with every vertex of $N_i(v)$.  

Note that $1\le \delta(G^*)\le\Delta(G^*)\leq 3$ and there is a natural bijection between the edges of $G$ and those of $G^*$.  Obviously, $G^*$ need not be connected. Let $H^*$ be a connected component of $G^*$. Now, we construct a proper and strong majority edge-coloring with 8 colors. 

Suppose first that all vertices of $H^*$ have degree 2 in $H^*$, that is, $H^*$ is a cycle $C_m$. It is easy to find a proper majority 4-edge-coloring in this case. Namely, take four colors, say 1,2,3,4, and color the consecutive edges of $C_m$ with three colors periodically: $1,2,3,1,2,3\ldots$. If $m\equiv 0\mod 3$, then we are done. If $m\equiv 1 \mod 3$, then we put color 4 on the last edge of $C_m$. Otherwise, for $m\equiv 2 \mod 3$, we color the last five  edges of $C_m$ with $1,4,2,3,4$.   
 
Let then $\Delta(H^*)=3$. By Vizing's theorem, $G^*$ has a proper edge-coloring $c$ with four colors. 

Let $e=xy$ be any edge of $H^*$. If both vertices $x,y$ are of odd degree, 1 or 3, then the edge $e$ is majority colored by $c$, for $c$ is proper. Consequently, if $e$ is not majority colored in $H^*$, then it is incident to a vertex of degree 2. 

Consider a maximal induced path $P=x_0,\ldots,x_r$ in $H^*$ with $r\ge 2$. Denote by $H^*_P$ the subgraph of $H^*$ induced by the vertices of $P$ and possible neighbors of the end-vertices $x_0$ and $x_r$. If $H^*_P$ contains an edge that is not majority colored, then we recolor some edges of $P$ as follows.

Suppose first that $P$ is a pendant path in $H^*$ and $d_{H^*}(x_r)=1$. If $d_{H^*}(x_0)=1$, then $H^*_P$ coincides with the path $P$. For any set $C'\subset C$ with $|C'|=3$, there exists a proper coloring $c':E(P)\rightarrow C'$ such that all  inner edges of $P$ are majority colored. 
Let then $d_{H^*}(x_0)=3$, and let $u_1,u_2$ be the two neighbors of $x_0$ outside of $P$. If $x_0u_1$ or $x_0u_2$ is not majority colored, then we recolor the edge $x_0x_1$ with a color from $C$ distinct from the colors of the six edges incident to $u_1$ and $u_2$. Then it is easy to see that we have at least $8-3=5$ free colors for the edge $x_1x_2$. Then we recolor subsequent edges $x_2x_3,\ldots, x_{r-1}x_r$ having 6 choices for each. It follows that any pair of free colors out of five can be used for the subpath $x_0x_1x_2$. We will use this fact later.

Now, assume that both end-vertices $x_0,x_r$ are of degree 3 in $H^*$. Denote by $v_1,v_2$ the two neighbors of $x_r$ outside $P$. We may need to recolor some edges of $P$ to obtain a proper coloring of $H^*_P$
such that every edge of $H^*_P$ is majority colored. 

Let $r=2$, and suppose that the edge $x_0x_1$ is not majority colored. It follows that the color $c(x_0x_1)$ also appears on one of the edges $x_0u_1$ or $x_0u_2$. To make the edge $x_0x_1$ majority colored, we recolor the edge $x_1x_2$. To do this, we cannot use the colors $c(x_0u_1), c(x_0u_2), c(x_0x_1),c(x_2v_1), c(x_2v_2)$. Moreover, if the vertices $v_1,v_2$ are of degree 2, then we cannot use the two colors of the edges outside $H^*_P$ incident to them. In total, we may have seven forbidden colors for the edge $x_1x_2$, so we may be forced to use an eighth color. By symmetry, the same arguments are used to show that we can recolor the edge $x_0x_1$ to make the edge $x_1x_2$ majority colored. 

Now, suppose that $r= 3$. If  $x_0x_1$ or $x_2x_3$ is not majority colored, then we can recolor the edge $x_1x_2$ with any color different from $c(x_0u_1), c(x_0u_2), c(x_0x_1)$ and from $c(x_2x_3),c(x_3v_1),c(x_3v_2)$. Then the subgraph $H^*_P$ gets a proper and strong majority coloring.  

Let then $r\ge 4$. If the edge $x_0x_1$ is not majority colored, we first recolor the edge $x_1x_2$. We choose  any color except $c(x_0u_1), c(x_0u_2), c(x_0x_1)$ for the edge $x_1x_2$. Analogously, if the edge $x_{r-1}x_r$ is not majority colored, we recolor the edge $x_{r-2}x_{r-1}$ with any color except  $c(x_{r-1}x_r),c(x_rv_1), c(x_rv_2)$. Then it is easy to see how to recolor the edges $x_2x_3,\ldots,x_{r-3}x_{r-2}$ with three colors from $C\setminus\{c(x_1x_2),c(x_{r-2}x_{r-1})\}$ to obtain a proper and strong majority coloring of $H^*_P$. 

Thus, we obtain a proper edge-coloring of every connected component of $G^*$, where every edge is majority colored except, possibly, some pendant edges of $G$. The sum of these colorings yields a strong majority edge-coloring of the graph $G$. This is because each coloring is proper, so a new component can add two or three edges with different colors incident to an end-vertex of a given edge. If an edge $e$ is not majority colored, then $e=x_0x_1$ belonged in $H^*$ to a pendant path $x_0,x_1,\ldots x_r$ of length $r\ge 2$. As we have noticed above, we can recolor the edges $x_0x_1, x_1x_2$ so that $e=x_0x_1$ becomes majority colored.
\end{proof}

\end{document}
